\documentclass[10pt]{article}
\usepackage{icmlko}

\icmlmeta{IP 파운데이션 월드모델}{65}{판정하는 코드에 규약을 심었다}{규칙은 사람이 어기고 구조는 안 어긴다 --- 노트 64의 처방을 모듈로 만든다}

\begin{document}

\twocolumn[
\icmlheader

\begin{abstract}
노트 64는 같은 종류의 오진을 셋 모으고 처방을 ``진단은 본 코드의 함수를 불러
쓴다''는 \emph{규칙}으로 적었다. 규칙은 급할 때 어긴다 --- 세 번 어긴 것이
바로 그 증거다. 구조로 옮겼다. \texttt{state/audit.py}에 판정에 필요한 것을
모았고, 핵심 설계는 하나다. \textbf{설정을 함수로 표현한다.} 배선 변경 · 축
끄기 · 대응 깨기를 전부 \texttt{(도메인, 축이름) $\to$ (도메인, 축이름)} 함수로
두면, 짝지은 붓스트랩이 \textbf{복원추출을 먼저 하고 설정을 뒤에 적용}할 수
있다. 노트 54의 버그 --- 복원추출 뒤에 고유 축을 붙여 레코드 대응이 깨지고
구간 부호가 뒤집힌 것 --- 이 이 구조에서는 \emph{만들 방법이 없다}. 레코드
색인도 여섯 디렉터리를 한꺼번에 훑게 해 노트 63의 미매칭 문제를 없앴다.
자기 점검을 붙였다 --- 이미 꺼져 있는 축을 끄면 $\Delta$가 정확히 0이어야 한다.
\textbf{0.0000이 나온다.} 나머지 둘도 앞선 노트의 값을 재현한다(팝업 매장축
$-$0.0077 대 노트 52의 $-$0.0074, 웹툰 대응 깨기 $-$0.0341 대 노트 53의
$-$0.0433). 앞으로 모든 판정이 이 모듈을 거친다.
\end{abstract}
]

\section{규칙에서 구조로}

노트 64의 처방은 규칙이었다. 그런데 그 노트가 모은 세 사례가 전부 ``급할 때
규약을 안 지킨'' 것이므로, 규칙을 하나 더 얹는 것은 같은 실패를 반복하는 길이다.

\begin{quote}\fontsize{9.5}{11.5}\selectfont
\textbf{설정을 함수로 표현한다.} \texttt{(도메인, 축이름) $\to$ (도메인, 축이름)}
꼴로 두면 짝지은 붓스트랩이 복원추출을 먼저 하고 설정을 뒤에 적용한다.
순서를 틀릴 방법이 없다.
\end{quote}

\begin{table}[t]
\centering\tsize
\caption{모듈이 담은 것.}
\begin{tabular}{ll}
\toprule
\texttt{domains()} & 정본 도메인 $+$ 축 이름 \\
\texttt{rho()} & 정본 판정치(노트 43) \\
\texttt{records()} & 여섯 디렉터리 레코드 색인 \\
\texttt{paired()} & 짝지은 붓스트랩 구간 \\
\texttt{cfg\_*()} & 설정 함수(끄기 · 섞기 · 배선) \\
\bottomrule
\end{tabular}
\end{table}

\texttt{records()}는 여섯 디렉터리를 한꺼번에 훑는다 --- 노트 63이 한 곳만 봐서
7건을 놓친 문제가 구조적으로 사라진다. 색인 1,316건이다.

\section{자기 점검}

도구를 새로 만들면 그 도구부터 검사해야 한다. 답을 아는 조건을 셋 넣었다.

\begin{figure}[t]
\centering
\includegraphics[width=\textwidth]{figs/audit.pdf}
\caption{자기 점검 셋. 첫 줄은 정확히 0이 나와야 한다.}
\end{figure}

\begin{table}[t]
\centering\tsize
\caption{자기 점검. 짝지은 붓스트랩 60회.}
\begin{tabular}{lrl}
\toprule
조건 & $\Delta\rho$ & 기대 \\
\midrule
게임 매장축 끄기(이미 꺼짐) & \textbf{$+$0.0000} & 정확히 0 \\
팝업 매장축 끄기 & $-$0.0077 & 노트 52: $-$0.0074 \\
웹툰 매장축 대응 깨기 & $-$0.0341 & 노트 53: $-$0.0433 \\
\bottomrule
\end{tabular}
\end{table}

첫 줄이 요점이다. 게임 매장 노출도는 노트 43에서 꺼 두었으므로 다시 꺼도 아무
일이 없어야 한다. 구간까지 $[+0.0000,\ +0.0000]$이다 --- 복제마다 두 설정이
완전히 같은 값을 낸다는 뜻이고, \textbf{설정 적용이 복원추출 뒤에 일어난다는
것의 직접 증거}다. 노트 54의 버그가 있었다면 여기서 0이 아닌 값이 나온다.

나머지 둘은 앞선 노트의 값을 재현한다. 웹툰 쪽이 $-$0.0341 대 $-$0.0433으로
조금 다른데, 노트 53은 씨앗 여럿의 평균이었고 여기서는 하나다.

\claimx{판정 도구에 자기 점검을 붙이면 진단 코드의 결함을 즉시 잡을 수 있다.
``이미 꺼진 축을 끄면 0''처럼 답을 아는 조건이 좋은 점검이다.}{점검 셋으로는
모든 결함을 못 잡는다. 특히 축 이름 정규화 같은 것은 이 셋에 안 걸린다.}

\section{앞으로}

이 모듈이 없던 예순네 노트의 판정은 각자 다른 임시 코드로 냈다. 그중 셋이
틀렸다는 것을 안다. 나머지는 모른다.

\begin{quote}\fontsize{9.5}{11.5}\selectfont
소급해서 전부 다시 재는 것은 하지 않는다 --- 계산이 크고, 틀린 것이 있다면
어차피 다음 노트에서 재현이 안 될 것이다. 대신 \textbf{앞으로의 판정은 전부
이 모듈을 거친다.}
\end{quote}

\section{한계}

\begin{itemize}[leftmargin=1.1em,itemsep=0.15em]
\item 자기 점검이 셋뿐이다. 답을 아는 조건을 더 만들어야 한다.
\item 배선 변경 설정 함수를 아직 안 넣었다. 지금은 축 끄기와 대응 깨기만 있다.
\item 앞선 노트의 판정을 소급 검증하지 않았다.
\item 모듈을 쓰라고 강제하는 장치가 없다. 여전히 임시 코드를 짤 수 있다.
\item 붓스트랩 회수를 60으로 낮춰 점검했다. 실제 판정은 200회 이상이어야 한다.
\end{itemize}

\section{다음}

\begin{enumerate}[leftmargin=1.3em,itemsep=0.15em]
\item 배선 변경 설정 함수를 모듈에 넣는다.
\item 자기 점검을 늘린다 --- 축을 두 번 끄면 한 번 끈 것과 같아야 한다 같은 것.
\item 웹툰 전체 수집 완료 후 전면 재측정(2,397/3,973 진행 중).
\item 앞선 노트 중 재현이 필요한 판정을 골라 이 모듈로 다시 잰다.
\end{enumerate}

\vspace{0.6em}
\hrule height 0.3pt
\vspace{0.5em}
{\footnotesize
\textbf{참고문헌}\\[0.3em]
Hoare, C. A. R. (1981). The emperor's old clothes. \emph{CACM}, 24(2), 75--83.\\[0.15em]
Efron, B., Tibshirani, R. J. (1993). \emph{An Introduction to the Bootstrap}.
Chapman \& Hall.\\[0.15em]
Bouthillier, X. et al. (2021). Accounting for variance in machine learning
benchmarks. \emph{MLSys}, 3, 747--769.\\[0.15em]
Spearman, C. (1904). The proof and measurement of association between two things.
\emph{American Journal of Psychology}, 15(1), 72--101.\\[0.15em]
Ben-David, S. et al. (2010). A theory of learning from different domains.
\emph{Machine Learning}, 79(1--2), 151--175.
}

\end{document}
